\documentclass[leqno]{jsarticle}
\usepackage{amssymb}
\usepackage{amsthm}
\usepackage{amsmath}
\usepackage{comment}
\usepackage[all]{xy}
%定理環境 thmはあまり使わずpropをなるべく使う。
%主張は基本的に証明の中で使い、そのため番号を付けない。
%注意環境を付けてみた。
\newtheorem{thm}{定理}
\newtheorem{prop}[thm]{命題}
\newtheorem{cor}[thm]{系}
\newtheorem{lem}[thm]{補題}
\newtheorem{df}[thm]{定義}
\newtheorem{exam}[thm]{例}
\newtheorem{fact}[thm]{事実}
\newtheorem*{claim}{主張}
\newtheorem*{cau}{注意}
\renewcommand{\proofname}{{\bf 証明}}
%矢印たち。
%\toは\rightarrowと同じ。\getsは\leftarrowと同じ。同値は\iff
%上向き矢はuparrow逆はdownarrow
\newcommand{\To}{\Longrightarrow}
\newcommand{\Gets}{\Longleftarrow}
%黒板太字
\newcommand{\nn}{\mathbb{N}}
\newcommand{\zz}{\mathbb{Z}}
\newcommand{\qq}{\mathbb{Q}}
\newcommand{\rr}{\mathbb{R}}
\newcommand{\cc}{\mathbb{C}}
%無限大
\newcommand{\mgn}{\infty}
%差集合は\setminusで統一する。等号の否定は\neq
%真部分集合は\subsetneqq　偏微分は\partial
%ディスプレイスタイル。\dfracでディスプレイ分数
\newcommand{\dpst}{\displaystyle}
\newcommand{\ep}{\varepsilon}

%空集合は\Oを使う！！！！！！！！！！！！
\newcommand{\supp}{\text{{\rm supp}}}
\newcommand{\wspp}{\text{{\rm wspp}}}

\title{正規性の特徴付け}
\author{一色三良(concious77)}
\date{\today}

\begin{document}
\maketitle
この文章では正規空間の様々な特徴付けを紹介する。ここでは正規空間とは$T_1$を仮定しない意味で用いる。

\begin{df}[細分]$\mathfrak{A},\mathfrak{B}を$集合$X$の部分集合の族とする。このとき$\mathfrak{A}$が$\mathfrak{B}$の細分であるとは
\[
\forall B\in \mathfrak{B}\ \exists A\in \mathfrak{A}\ A\subset B
\]
となることである。
\end{df}

\begin{df}[点有限族]集合$X$の部分集合の族$\mathfrak{A}$が点有限であるとは$X$の任意の点$x$に対して$x\in A$となる$A\in \mathfrak{A}$が有限個しかないことである。

\end{df}

\begin{df}[局所有限族]位相空間$X$に於いて$X$の部分集合の族$\mathfrak{A}$が局所有限であるとは$X$の各点$x$について$x$の近傍$N$が存在して
\[
N\cap A\neq \emptyset となるA\in \mathfrak{A}は有限個
\]
を充たすことである。
\end{df}

\begin{df}[開収縮]位相空間$X$の開被覆$\mathcal{U}$に対して$\{V_U\ |\ U\in \mathcal{U}\}$が開収縮であるとは$\{V_U\ |\ U\in \mathcal{U}\}$が$X$の開被覆であって、任意の$U\in \mathcal{U}$について
\[
CL(V_U)\subset U
\]
を充たすこと。
\end{df}

\begin{df}[閉収縮]位相空間$X$の開被覆$\mathcal{U}$に対して$\{F_U\ |\ U\in \mathcal{U}\}$が閉収縮であるとは$\{F_U\ |\ U\in \mathcal{U}\}$が$X$の閉被覆であって、任意の$U\in \mathcal{U}$について
\[
F_{U}\subset U
\]
を充たすこと。
\end{df}


\begin{df}[Weak supportとSupport]
関数$fX\to [0,1]$の{\bf 弱台}$\wspp (f)$を
\[
\wspp (f)=\{x\in X\ |\ f(x)>0\}
\]
と定義する。\par
関数$fX\to [0,1]$の{\bf 台}$\supp (f)$を
\[
\supp (f)=CL(\supp(f))=CL(\{x\in X\ |\ f(x)>0\})
\]

\end{df}

以下様々な種類の単位の分割を定義する。
%弱い単位の分割
\begin{df}[弱い単位の分割]位相空間$X$上の実数値連続関数の族$\{f_i\}_{i\in I}$が$X$の開被覆$\mathcal{U}$に従属する弱い単位の分割であるとは$\{\wspp f_i\}_{i\in I}$が$\mathcal{U}$の細分であり
\begin{align*}
&\sum_{i\in I} f_i=1
\end{align*}
を充たすこと
\end{df}
%強い単位の分割
\begin{df}[強い単位の分割]位相空間$X$上の実数値連続関数の族$\{f_i\}_{i\in I}$が$X$の開被覆$\mathcal{U}$に従属する強い単位の分割であるとは$\{\supp f_i\}_{i\in I}$が$\mathcal{U}$の細分であり
\begin{align*}
&\sum_{i\in I} f_i=1
\end{align*}
を充たすこと
\end{df}
%点有限な弱い単位の分割
\begin{df}[点有限な弱い単位の分割]位相空間$X$上の実数値連続関数の族$\{f_i\}_{i\in I}$が$X$の開被覆$\mathcal{U}$に従属する弱い単位の分割であるとは$\{\wspp f_i\}_{i\in I}$が$\mathcal{U}$の細分であり
\begin{align*}
&\sum_{i\in I} f_i=1\\
&\text{$\{\wspp f_i\}_{i\in I}$は点有限}
\end{align*}
を充たすこと
\end{df}

%点有限な強い単位の分割
\begin{df}[点有限な弱い単位の分割]位相空間$X$上の実数値連続関数の族$\{f_i\}_{i\in I}$が$X$の開被覆$\mathcal{U}$に従属す強い単位の分割であるとは$\{\supp f_i\}_{i\in I}$が$\mathcal{U}$の細分であり
\begin{align*}
&\sum_{i\in I} f_i=1\\
&\text{$\{\supp f_i\}_{i\in I}$は点有限}
\end{align*}
を充たすこと
\end{df}

%局所有限な弱い単位の分割
\begin{df}[局所有限な弱い単位の分割]位相空間$X$上の実数値連続関数の族$\{f_i\}_{i\in I}$が$X$の開被覆$\mathcal{U}$に従属する局所有限な弱い単位の分割であるとは$\{\wspp f_i\}_{i\in I}$が$\mathcal{U}$の細分であり
\begin{align*}
&\sum_{i\in I} f_i=1\\ 
&\text{$\{\wspp f_{i}\}_{i\in I}$が局所有限}
\end{align*}
を充たすこと
\end{df}

%局所有限な強い単位の分割
\begin{df}[局所有限な強い単位の分割]位相空間$X$上の実数値連続関数の族$\{f_i\}_{i\in I}$が$X$の開被覆$\mathcal{U}$に従属する局所有限な強い単位の分割であるとは$\{\supp f_i\}_{i\in I}$が$\mathcal{U}$の細分であり
\begin{align*}
&\sum_{i\in I} f_i=1\\ 
&\text{$\{\supp f_{i}\}_{i\in I}$が局所有限}
\end{align*}
を充たすこと
\end{df}

\section*{関数の拡張}
\begin{df}
位相空間$X$について、$X$の任意の交わらない２つの閉集合の組$E,F$について連続関数$f:X\to [0,1]$が存在して
\[
f|E\equiv 0,\ f|F\equiv 1
\]
が成り立つとき、$X$は条件$(U)$を充たすという。
\end{df}


\begin{df}
位相空間$X$について、任意の$X$の閉集合$A$と$A$上の連続関数$f:A\to [0,1]$に対して、連続関数$g:X\to [0,1]$が存在して
\[
g|A=f
\]
を充たすとき$X$は条件$(T)$を充たすという。
\end{df}

\begin{df}
位相空間$X$について、任意の$X$の閉集合$A$と$A$上の連続関数$f:A\to \rr$に対して、連続関数$g:X\to \rr$が存在して
\[
g|A=f
\]
を充たすとき$X$は条件$(R)$を充たすという。
\end{df}

\begin{thm}
位相空間$X$について以下は同値である。
\begin{align}
&\mbox{正規性}\\ 
&\mbox{条件$(U)$を充たす}\\ 
&\mbox{条件$(T)$を充たす}\\
&\mbox{条件$(R)$を充たす}
\end{align}
\end{thm}
\section*{$\aleph_0$-族正規性}
\begin{df}[$\aleph_0$-族正規性もしくは可算族正規性]
位相空間$X$が$\aleph_0$-族正規であるとは$X$の閉集合からなる任意の{\bf 可算で}疎な族$\mathcal{A}$に対して$A\subset U_A$となる開集合の族$\{U_{A}\ |\ A\in \mathcal{A}\}$が存在して
\[
A\neq B\To U_A\cap U_B=\O
\]
が成り立つことである。\footnote{$T_1$は要請しないことにする}明らかに$\aleph_0$-族正規空間は正規である。

\end{df}
\begin{thm}
位相空間$X$に於いて正規性と$\aleph_0$-族正規性は同値
\end{thm}
\begin{proof}
正規空間が$\aleph_0$-族正規であることを示せばいい。
\end{proof}

\section*{縮小}
\setcounter{equation}{0}
\begin{thm}
位相空間$X$について以下は同値である。
\begin{align}
&\mbox{正規性}\\ 
&\mbox{点有限開被覆に開収縮が存在する}\\ 
&\mbox{局所有限開被覆に開収縮が存在する}\\ 
&\mbox{有限開被覆に開収縮が存在する}\\ 
&\mbox{点有限開被覆に閉収縮が存在する}\\
&\mbox{局所有限開被覆に閉収縮が存在する}\\  
&\mbox{有限開被覆に閉収縮が存在する}
\end{align}
\end{thm}
\begin{proof}
$(1)\To (2)$の証明は省略する。これは所謂ウリゾーンの補題である。\par
交わらない２つの閉集合$E,F$について$f|E\equiv 0,\ f|F\equiv 1$となる連続関数$f:\to [0,1]$について、$U=f^{-1}([0,1/2)),\ V=f^{-1}((1/2,1])$を考えれば$U,V$は互いに交わらない開集合で$E\subset U,\ F\subset V$なので結局$(2)\To (1)$が成り立つ。また$(3)$を仮定したとき互いに交わらない２つの閉集合$E,F$について関数$f:E\cup F\to [0,1]$を$E$上で$0$、$F$上で$1$の値を取る関数とすればこれは連続関数で$(3)$を用いれば連続関数$g:X\to [0,1]$で$g|E\equiv 0,\ g|F\equiv 1$となる関数が得られるので$(2)$が成り立つ。以上の簡単な考察から
\[
(1)\iff (2)\Gets (3)
\]
が分かった。また、$(1)\To (3)$はティーチェの拡張定理としてよく知られているので証明を省略する。そして
以下の論理包含図式は明らかであろう。
\[\xymatrix{
& (4) \ar@{=>}[r]\ar@{=>}[d]& (5)\ar@{=>}[r]\ar@{=>}[d] & (6)\ar@{=>}[d] &\\ 
& (7)\ar@{=>}[r] & (8) \ar@{=>}[r]& (9) & 
}\]
これらを踏まえて以下では
\[
(1)\To (4),\ (9)\To (1),\ (5)\To(10),\ (10)\To (2)
\]
を証明する。\par
[$(1)\To(4)$の証明]\par
$\{U_{i}\}_{i\in I}$を$X$の任意の点有限開被覆とし、$I$は整列しているとする。超限帰納法で
\begin{align*}
&(hoge1)\ CL(V_i)\subset U_i\\ 
&(hoge2)\ \forall k\in I\ \{V_{i}\}_{i\le k}\cup\{U_i\}_{k<i}\mbox{が$X$の被覆を成している}
\end{align*}
となる$V_i$を構成する。$0$を$I$の最小元として、まず$V_0$を構成する。$F_0=X\setminus \bigcup_{k\neq0}U_k$と置くとこれは閉集合で$\{U_i\}_{i\in I}$が被覆を成すことから
\[
F_0\subset U_0
\]
となる。ここで$X$の正規性を用いると
\[
F_0\subset V\subset CL(V)\subset U_0
\]
となる開集合$V$が存在する。このような$V$の一つを$V_0$とおくと
\[
CL(V_0)\subset U_0
\]
を明らかに充たし
更に$\{V_0\}\cup\{U_i\}_{i>0}$が$X$の被覆を成している。\par
次に$i<k$となる任意の$i$については$(a),(b)$を充たす$V_i$が構成できているとしよう。このとき
\[
F_k=X\setminus \left(\bigcup_{i<k}V_i\cup \bigcup_{k<i}U_i\right)
\]
と定義する。さて
\[
F_k\subset U_k
\]
となることを見よう。もし$x\in F_k(X\setminus U_k)$が存在したとしよう。$\{U_i\}_{i\in I}$は点有限なので$x\in U_{a}$となる$a\in I$は有限個であり、このようなもの全体を$\{a_1,a_2,\cdots,a_n\}$としよう。そして$b=\max\{a_1,a_2\cdots ,a_n\}$と置くと$F_k$の定義と$x\in F_k$と$x\not\in U_k$より$b<k$であり、$b<i$となる$i$について$x\not\in U_i$であるので$b$に対する$(hoge2)$から$x\in \bigcup_{i\le b}V_b$であるが、$\bigcup_{i\le b}V_b\subset \bigcup_{i<k}V_i$と$F_k$の定義から$x\not\in F_k$になり、矛盾する。よって$F_k\subset U_k$\par
さてさてここで$X$の正規性から
\[
F_k\subset V\subset CL(V)\subset U_k
\]
となる開集合$V$が存在するが、このような$V$のひとつを$V_k$としよう。このとき、$(hoge1),\ (hoge2)$が充たされることは見やすい。\par
最後にこのように作られた$\{V_i\}_{i\in I}$が$X$の被覆になることを見よう。$x\in X$を任意に与える
$\{U_i\}_{i\in I}$は点有限なので$x\in U_{a}$となる$a\in I$は有限個であり、このようなもの全体を$\{a_1,a_2,\cdots,a_n\}$としよう。そして$b=\max\{a_1,a_2\cdots ,a_n\}$と置くと$b<i$ならば$x\not\in U_i$なので$(hoge2)$から$x\in V_{j}$となる$V_j$が存在する。よって$\{V_i\}_{\in I}$は被覆となる。
\par [$(1)\To(4)$の証明終わり]\par
[$(9)\To(1)$の証明]\par
$E,F$を交わらない２つの$X$の閉集合としよう。このとき$\{X\setminus E,X\setminus F\}$は$X$の有限開被覆なので$(9)$から閉収縮$\{A,B\}$が存在する。つまり$\{A,B\}$は
\[
A\subset X\setminus E,\ B\subset X\setminus F
\]
となる２つの閉集合である。このとき$U=X\setminus A,\ V=X\setminus B$と置くと、これらは開集合であり、更に$E\subset U,\ F\subset V$である。また$A\cup B=X$から
\[
U\cap V=X\setminus (A\cup B)=X\setminus X=\O
\]
なので結局$X$は正規である。
\par [$(9)\To(1)$の証明終わり]\par
[$(5)\To(10)$の証明]\par
$\{U_i\}_{i\in I}$を局所有限開被覆としよう。このとき$(5)$から$$\{U_i\}_{i\in I}$$の開収縮$\{V_i\}_{i\in I}$が存在する。またこれも局所有限である。さて$(1)$から$(9)$の同値性はもう知っている。$\{V_i\}_{i\in I}$に$(8)$を用いてこれの閉収縮$\|F_i\}_{i\in I}$が存在する。そして$F_i\subset V_i$なのでウリゾーンの補題から各$i\in I$毎に
\[
g_{i}|F_i\equiv 1,\ g_i|X\setminus V_i\equiv 0
\]
を充たす連続関数$g_i:X\to [0,1]$が存在する。ここで$CL(V_i)\subset U_i$なので
\[
\supp g_i\subset U_i
\]
であるから$\{\supp g_i\}_{i\in I}$は$\{U_i\}_{i\in I}$の細分である。
さて、局所有限性から連続関数$h:X\to[0,+\mgn)$を
\[
h(x)=\sum_{i\in I}g_i(x)
\]
と定義でき、更に$\{F_i\}_{i\in I}$が被覆を成していることから$g(x)>0$である。そして
\[
f_i(x)=\frac{g_i(x)}{h(x)}
\]
と連続関数$f_iX\to [0,1]$を定義する事ができて、これが$\{U_i\}_{i\in I}$に従属する局所有限な単位の分割であることは容易にわかる。

\par [$(5)\To(10)$の証明終わり]\par
[$(10)\To(2)$の証明]\par
$E,F$を交わらない２つの$X$の閉集合としよう。このとき$\mathcal{A}=\{X\setminus E,X\setminus F\}$は$X$の局所有限開被覆なのでこれに従属する局所有限な単位の分割$\{f_i\}_{i\in I}$が存在する。ここで$S\subset I$を
\[
\supp f_i\subset X\setminus E
\]
を充たす$\in I$全体の集合とし、
\[
g(x)=\sum_{s\in S}f_s(x)
\]
と定義すると局所有限性から$g$は連続関数であり、$s\in S$ならば$\supp f_s\subset X\setminus E$なので$g|E\equiv 0$である。また$\{supp f_i\}_{i\in I}$は$\mathcal{A}$の細分であるので$\supp f_i\cap E\neq\O$と$\supp f_{i}\cap F\neq\O$は同時に絶対に起らない。よって$x\in F$に対して$x\in \supp f_a$となる$f_a$を考えると必ず$\supp f_a\subset X\setminus E$なので$\sum_{i\in I} f_i=1$から
\[
g|F\equiv 1
\]
がわかる。よって$(2)$が成立する、
\par [$(10)\To(2)$の証明終わり]\par
\end{proof}

\section*{単位の分割}
\setcounter{equation}{0}
\begin{lem}
位相空間$X$とその開被覆$\mathcal{U}$について以下は同値
\begin{align}
&\mbox{$\mathcal{U}$に従属する弱い単位の分割が存在する}\\
&\mbox{$\mathcal{U}$に従属する局所有限な弱い単位の分割が存在する}
\end{align}
\end{lem}
\setcounter{equation}{0}
\begin{thm}[局所有限被覆と単位の分割]
位相空間$X$について以下は同値である。
\begin{align}
&\mbox{正規性}\\
&\mbox{任意の局所有限開被覆に対してそれに従属する弱い単位の分割が存在する}\\ 
&\mbox{任意の局所有限開被覆に対してそれに従属する強い単位の分割が存在する}\\ 
&\mbox{任意の局所有限開被覆に対してそれに従属する点有限な弱い単位の分割が存在する}\\ 
&\mbox{任意の局所有限開被覆に対してそれに従属する点有限な強い単位の分割が存在する}\\ 
&\mbox{任意の局所有限開被覆に対してそれに従属する局所有限な弱い単位の分割が存在する}\\ 
&\mbox{任意の局所有限開被覆に対してそれに従属する局所有限な強い単位の分割が存在する}
\end{align}
\end{thm}
\setcounter{equation}{0}
\begin{thm}[有限被覆と単位の分割]
位相空間$X$について以下は同値である。
\begin{align}
&\mbox{正規性}\\
&\mbox{任意の有限開被覆に対してそれに従属する弱い単位の分割が存在する}\\ 
&\mbox{任意の有限開被覆に対してそれに従属する強い単位の分割が存在する}\\ 
&\mbox{任意の有限開被覆に対してそれに従属する点有限な弱い単位の分割が存在する}\\ 
&\mbox{任意の有限開被覆に対してそれに従属する点有限な強い単位の分割が存在する}\\ 
&\mbox{任意の有限開被覆に対してそれに従属する局所有限な弱い単位の分割が存在する}\\ 
&\mbox{任意の有限開被覆に対してそれに従属する局所有限な強い単位の分割が存在する}
\end{align}
\end{thm}
\setcounter{equation}{0}
%以下は成り立たないので修正
\begin{thm}[点有限被覆と単位の分割]
位相空間$X$について以下は同値である。
\begin{align}
&\mbox{正規性}\\ 
&\mbox{任意の点有限開被覆に対してそれに従属する弱い単位の分割が存在する}\\ 
&\mbox{任意の点有限開被覆に対してそれに従属する強い単位の分割が存在する}\\ 
&\mbox{任意の点有限開被覆に対してそれに従属する点有限な弱い単位の分割が存在する}\\ 
&\mbox{任意の点有限開被覆に対してそれに従属する点有限な強い単位の分割が存在する}\\ 
&\mbox{任意の点有限開被覆に対してそれに従属する局所有限な弱い単位の分割が存在する}\\ 
&\mbox{任意の点有限開被覆に対してそれに従属する局所有限な強い単位の分割が存在する}
\end{align}
\end{thm}

\section*{$\Delta$細分}
\setcounter{equation}{0}
\begin{thm}位相空間$X$について以下は正規性と同値
\begin{align}
\mbox{任意の開集合二枚からなる開被覆について$\Delta$細分となる有限開被覆が存在する。}\\
\mbox{任意の有限開被覆について$\Delta$細分となる有限開被覆が存在する。}\\
\mbox{任意の有限開被覆について$\Delta$細分となる開被覆が存在する。}\\
\mbox{任意の有限開被覆は正規被覆となる。}\\
\mbox{任意の局所有限開被覆について$\Delta$細分となる開被覆が存在する。}\\
\end{align}
\end{thm}

\begin{thebibliography}{9}
\bibitem{1}森田紀一,位相空間論,岩波全書331,岩波書店1981
\end{thebibliography}

\end{document}








